\section{Datenmodell und Zustandsverwaltung}
\label{sec:datenmodell}

\subsection{Task-Datenmodell}

Jeder Task im System wird durch die Datenklasse \texttt{Task}
(in \texttt{src/dispatcher/task.py}) repräsentiert.
Das Modell entspricht dem in der Aufgabenstellung geforderten
\texttt{task\_t}-Typ und enthält alle Pflichtfelder:

\begin{table}[H]
\centering
\begin{tabular}{llll}
\toprule
\textbf{Feld} & \textbf{Typ} & \textbf{Bedeutung} & \textbf{Länge} \\
\midrule
\texttt{task\_id}              & \texttt{str}        & Eindeutige Task-Nummer (int32-kompatibler Zähler) & -- \\
\texttt{task\_type}            & \texttt{str}        & Aufgabentyp (z.\,B. \texttt{sum}, \texttt{reverse}) & max. 32 \\
\texttt{payload}               & \texttt{str}        & Eingabedaten                                       & max. 1024 \\
\texttt{result}                & \texttt{str}        & Berechnetes Ergebnis (leer bis COMPLETED)           & max. 1024 \\
\texttt{status}                & \texttt{TaskState}  & Aktueller Zustand (Enum)                           & -- \\
\texttt{timestamp\_created}    & \texttt{int}        & Unix-Timestamp Erstellungszeitpunkt                & -- \\
\texttt{timestamp\_dispatched} & \texttt{int}        & Unix-Timestamp erster Dispatch-Zeitpunkt           & -- \\
\texttt{timestamp\_completed}  & \texttt{int}        & Unix-Timestamp Abschlusszeitpunkt                  & -- \\
\texttt{retry\_count}          & \texttt{int}        & Anzahl bisheriger Retry-Versuche                   & -- \\
\texttt{assigned\_worker}      & \texttt{str}        & Worker-ID des zugeordneten Workers                 & max. 64 \\
\bottomrule
\end{tabular}
\caption{Task-Datenfelder}
\end{table}

Die Längenbeschränkungen (\texttt{MAX\_TYPE\_LEN = 32}, \texttt{MAX\_PAYLOAD\_LEN = 1024},
\texttt{MAX\_WORKER\_LEN = 64}) werden im \texttt{PostTask}-Handler geprüft.
Überschreitungen werden mit gRPC-Fehlercode \texttt{INVALID\_ARGUMENT} abgelehnt.

\subsection{Task-Identifikation}

Task-IDs werden durch einen thread-sicheren, monoton steigenden Zähler
in \texttt{src/common/protocol.py} erzeugt:

\begin{lstlisting}[language=Python,caption={Thread-sichere Task-ID-Generierung}]
_task_id_lock = threading.Lock()
_task_id_counter = 0

def new_task_id() -> str:
    global _task_id_counter
    with _task_id_lock:
        _task_id_counter += 1
        return str(_task_id_counter)
\end{lstlisting}

Der Zähler wird als String gespeichert, ist aber int32-kompatibel --
notwendig für die Proto-Kompatibilität mit dem Worker-Code.

\subsection{Zustände der Zustandsmaschine}

Das System implementiert eine Zustandsmaschine mit \textbf{8 Zuständen}
gemäß Aufgabenstellung §5. Zwei davon sind Terminalzustände:

\begin{table}[H]
\centering
\begin{tabular}{lll}
\toprule
\textbf{Zustand} & \textbf{Bedeutung} & \textbf{Terminal?} \\
\midrule
\texttt{CREATED}    & Task angelegt, noch nicht eingereiht        & nein \\
\texttt{QUEUED}     & In Warteschlange, wartet auf Dispatch        & nein \\
\texttt{DISPATCHED} & An Worker gesendet, wartet auf Akzeptanz     & nein \\
\texttt{PROCESSING} & Vom Worker akzeptiert, wird verarbeitet      & nein \\
\texttt{TIMEOUT}    & Kein Ergebnis in Timeout-Zeitfenster         & nein \\
\texttt{RETRYING}   & Wird erneut eingeplant (Retry läuft)         & nein \\
\texttt{COMPLETED}  & Ergebnis erfolgreich empfangen               & \textbf{ja} \\
\texttt{FAILED}     & Verarbeitung endgültig fehlgeschlagen        & \textbf{ja} \\
\bottomrule
\end{tabular}
\caption{Task-Zustände und ihre Bedeutung}
\end{table}

\subsection{Erlaubte Zustandsübergänge}

Die Zustandsmaschine ist in \texttt{src/dispatcher/state\_machine.py}
implementiert. Ein Dictionary \texttt{VALID\_TRANSITIONS} definiert
alle erlaubten Übergänge:

\begin{lstlisting}[language=Python,caption={Erlaubte Zustandsübergänge (state\_machine.py)}]
VALID_TRANSITIONS: dict[TaskState, set[TaskState]] = {
    TaskState.CREATED:    {TaskState.QUEUED},
    TaskState.QUEUED:     {TaskState.DISPATCHED, TaskState.FAILED},
    TaskState.DISPATCHED: {TaskState.PROCESSING, TaskState.TIMEOUT,
                           TaskState.FAILED},
    TaskState.PROCESSING: {TaskState.COMPLETED,  TaskState.FAILED,
                           TaskState.TIMEOUT},
    TaskState.TIMEOUT:    {TaskState.RETRYING},
    TaskState.RETRYING:   {TaskState.DISPATCHED,  TaskState.FAILED},
    TaskState.COMPLETED:  set(),   # terminal
    TaskState.FAILED:     set(),   # terminal
}
\end{lstlisting}

Jeder Übergangsversuch der nicht in \texttt{VALID\_TRANSITIONS} gelistet ist,
wirft eine \texttt{InvalidTransitionError}-Exception.
Dies verhindert inkonsistente Zustände im Concurrent-Betrieb.

\subsection{Zustandsdiagramm}

\begin{figure}[H]
\centering
\begin{lstlisting}[language={},frame=single,basicstyle=\ttfamily\small]
                    PostTask
  [Client]  --------> CREATED --> QUEUED
                                     |
                        dispatch     |  (kein Worker)
                         Loop        |  re-enqueue
                             |       |
                             v       v
                         DISPATCHED <-------- RETRYING
                          |    |                ^   |
              ACK ok=True |    | Timeout        |   | retry_count >= MAX_RETRIES
                          v    v                |   v
                     PROCESSING --> TIMEOUT --->+  FAILED
                          |
               ReturnResult|  COMPLETED/FAILED
                          v
               COMPLETED (terminal) / FAILED (terminal)
\end{lstlisting}
\caption{Zustandsübergangsdiagramm (vereinfacht)}
\end{figure}

\subsection{Zustandsübergang-Implementierung}

Die Funktion \texttt{transition()} wendet einen Zustandsübergang an
und setzt automatisch die Timestamps:

\begin{lstlisting}[language=Python,caption={transition()-Funktion (state\_machine.py)}]
def transition(task: Task, new_state: TaskState) -> Task:
    allowed = VALID_TRANSITIONS.get(task.status, set())
    if new_state not in allowed:
        raise InvalidTransitionError(
            f"Task {task.task_id}: Uebergang {task.status} -> "
            f"{new_state} nicht erlaubt."
        )
    task.status = new_state

    if new_state == TaskState.DISPATCHED:
        task.timestamp_dispatched = int(time.time())
    elif new_state in (TaskState.COMPLETED, TaskState.FAILED):
        task.timestamp_completed = int(time.time())

    return task
\end{lstlisting}

\subsection{Thread-sichere Datenhaltung}

\subsubsection{TaskStore (Issue \#21)}

Der \texttt{TaskStore} hält alle Tasks in einem Dictionary im Arbeitsspeicher.
Alle Methoden (Add, Get, Update, AllTasks) sind unter einem einzigen
\texttt{threading.Lock()} serialisiert:

\begin{itemize}
  \item \texttt{add(task)}: Fügt neuen Task ein.
  \item \texttt{get(task\_id)}: Gibt Task per ID zurück (None wenn nicht gefunden).
  \item \texttt{update(task)}: Überschreibt bestehenden Task; wirft
        \texttt{KeyError} wenn nicht vorhanden.
  \item \texttt{all\_tasks()}: Gibt Snapshot aller Tasks zurück.
  \item \texttt{count\_by\_status(status)}: Zählt Tasks in einem Zustand.
  \item \texttt{worker\_load\_map()}: Gibt Anzahl aktiver Tasks pro Worker zurück.
\end{itemize}

\subsubsection{TaskQueue (Issue \#21)}

Die \texttt{TaskQueue} kapselt \texttt{queue.Queue} aus der Python-Standardbibliothek.
Diese ist von Haus aus thread-sicher (FIFO, blockierendes \texttt{get()}).
Die Kapselung ermöglicht, künftig eine Priorisierung einzuführen,
ohne den Rest des Dispatchers anzupassen.

\begin{lstlisting}[language=Python,caption={TaskQueue-Implementierung (task\_queue.py)}]
class TaskQueue:
    def __init__(self) -> None:
        self._q: queue.Queue[Task] = queue.Queue()

    def enqueue(self, task: Task) -> None:
        self._q.put(task)

    def dequeue(self, timeout: float = 1.0) -> Optional[Task]:
        try:
            return self._q.get(timeout=timeout)
        except queue.Empty:
            return None
\end{lstlisting}

\subsection{Nebenläufigkeit im Dispatcher (Issue \#21)}

Der Dispatcher verarbeitet mehrere Tasks gleichzeitig durch folgendes Design:

\begin{enumerate}
  \item Der \texttt{DispatchLoop}-Thread liest sequenziell aus der Queue
        (kein Race auf \texttt{Queue.get()}).
  \item Jeder Dispatch-Vorgang wird in einem eigenen Thread des
        \texttt{ThreadPoolExecutor(max\_workers=10)} ausgeführt.
  \item Dadurch können bis zu 10 Tasks gleichzeitig \textit{in-flight} sein
        (an Worker gesendet, warten auf Ergebnis).
  \item Alle Zugriffe auf \texttt{TaskStore} sind durch dessen internen Lock
        serialisiert -- keine Race Conditions beim Zustandswechsel.
  \item Timeout-Timer laufen ebenfalls in separaten Daemon-Threads.
        Der Zugriff auf das Timer-Dictionary \texttt{\_timers} ist
        durch \texttt{\_timers\_lock} geschützt.
\end{enumerate}
